\renewcommand{\bibname}{Bibliographic Guide}

\begin{thebibliography}{99}
\addcontentsline{toc}{chapter}{Bibliographic Guide}

\item[]
\textbf{How this bibliography is used.}
The references below are source traces for theorem-boundary statements,
historical origin of constructions, and convention checks.  A citation is not
a proof in this monograph.  Every result used in the logical development must
still be derived in the local text, stated with its hypotheses as a
quoted theorem, or marked as a hypothesis, controlled approximation,
conjecture, or open problem.

\item[]
This first seed covers the most frequently quoted foundational lineages:
Wightman and OS reconstruction
\cite{Wightman1956,StreaterWightman1964,Jost1965,OS1973,OS1975,
GubinelliHofmanova2021,KravchukQiaoRychkov2021}, scattering theory
\cite{Haag1958,Ruelle1962,LSZ1955}, perturbative and
Wilsonian renormalization \cite{BPH1957,Hepp1966,Zimmermann1969,Polchinski1984},
AQFT superselection and modular theory
\cite{HaagKastler1964,ReehSchlieder1961,DHR1971,DHR1974,DR1989,DR1990,
GelfandNaimark1943,Segal1947,Takesaki1970,Connes1973,DAntoniLongo1983,
Borchers1992},
constructive and stochastic routes \cite{GlimmJaffe1987,Hairer2014}, CFT and
VOA structure \cite{Segal2004,KontsevichSegal2021,Zhu1996,Huang2005,DO1994,ZZ1996}, anomalies and index
theory
\cite{Seeley1967,Gilkey1984,AtiyahSinger1968,APS1975I,APS1975II,APS1976III,
BismutFreed1986,DaiFreed1994,Witten1982}, instantons
\cite{tHooft1976Pseudoparticle}, supersymmetric dynamics
\cite{SeibergWitten1994I,SeibergWitten1994II,Seiberg1994}, and locality
laboratories \cite{LiebRobinson1972}.  Later passes should attach additional
citations at the precise theorem-boundary sites and expand this guide.

\bibitem{Wightman1956}
A.~S. Wightman.
Quantum field theory in terms of vacuum expectation values.
\emph{Physical Review} \textbf{101} (1956), 860--866.

\bibitem{StreaterWightman1964}
R.~F. Streater and A.~S. Wightman.
\emph{PCT, Spin and Statistics, and All That}.
Benjamin, 1964.

\bibitem{Jost1965}
R.~Jost.
\emph{The General Theory of Quantized Fields}.
American Mathematical Society, 1965.

\bibitem{Haag1958}
R.~Haag.
Quantum field theories with composite particles and asymptotic conditions.
\emph{Physical Review} \textbf{112} (1958), 669--673.

\bibitem{Ruelle1962}
D.~Ruelle.
On the asymptotic condition in quantum field theory.
\emph{Helvetica Physica Acta} \textbf{35} (1962), 147--163.

\bibitem{LSZ1955}
H.~Lehmann, K.~Symanzik, and W.~Zimmermann.
On the formulation of quantized field theories.
\emph{Il Nuovo Cimento} \textbf{1} (1955), 205--225.

\bibitem{OS1973}
K.~Osterwalder and R.~Schrader.
Axioms for Euclidean Green's functions.
\emph{Communications in Mathematical Physics} \textbf{31} (1973), 83--112.

\bibitem{OS1975}
K.~Osterwalder and R.~Schrader.
Axioms for Euclidean Green's functions II.
\emph{Communications in Mathematical Physics} \textbf{42} (1975), 281--305.

\bibitem{GubinelliHofmanova2021}
M.~Gubinelli and M.~Hofmanov\'a.
A PDE construction of the Euclidean \(\Phi^4_3\) quantum field theory.
\emph{Communications in Mathematical Physics} \textbf{384} (2021), 1--75.

\bibitem{KravchukQiaoRychkov2021}
P.~Kravchuk, J.~Qiao, and S.~Rychkov.
Distributions in conformal field theory II: Minkowski space.
\emph{Journal of High Energy Physics} \textbf{08} (2021), 094.

\bibitem{BPH1957}
N.~N. Bogoliubov and O.~S. Parasiuk.
On the multiplication of causal functions in the quantum theory of fields.
\emph{Acta Mathematica} \textbf{97} (1957), 227--266.

\bibitem{Hepp1966}
K.~Hepp.
Proof of the Bogoliubov--Parasiuk theorem on renormalization.
\emph{Communications in Mathematical Physics} \textbf{2} (1966), 301--326.

\bibitem{Zimmermann1969}
W.~Zimmermann.
Convergence of Bogoliubov's method of renormalization in momentum space.
\emph{Communications in Mathematical Physics} \textbf{15} (1969), 208--234.

\bibitem{Polchinski1984}
J.~Polchinski.
Renormalization and effective Lagrangians.
\emph{Nuclear Physics B} \textbf{231} (1984), 269--295.

\bibitem{HaagKastler1964}
R.~Haag and D.~Kastler.
An algebraic approach to quantum field theory.
\emph{Journal of Mathematical Physics} \textbf{5} (1964), 848--861.

\bibitem{ReehSchlieder1961}
H.~Reeh and S.~Schlieder.
Bemerkungen zur unit\"ar\"aquivalenz von Lorentz-invarianten Feldern.
\emph{Nuovo Cimento} \textbf{22} (1961), 1051--1068.

\bibitem{GelfandNaimark1943}
I.~M. Gelfand and M.~A. Naimark.
On the imbedding of normed rings into the ring of operators in Hilbert space.
\emph{Matematicheskii Sbornik} \textbf{12(54)} (1943), 197--217.

\bibitem{Segal1947}
I.~E. Segal.
Irreducible representations of operator algebras.
\emph{Bulletin of the American Mathematical Society} \textbf{53} (1947),
73--88.

\bibitem{DHR1971}
S.~Doplicher, R.~Haag, and J.~E.~Roberts.
Local observables and particle statistics I.
\emph{Communications in Mathematical Physics} \textbf{23} (1971), 199--230.

\bibitem{DHR1974}
S.~Doplicher, R.~Haag, and J.~E.~Roberts.
Local observables and particle statistics II.
\emph{Communications in Mathematical Physics} \textbf{35} (1974), 49--85.

\bibitem{DR1989}
S.~Doplicher and J.~E.~Roberts.
A new duality theory for compact groups.
\emph{Inventiones Mathematicae} \textbf{98} (1989), 157--218.

\bibitem{DR1990}
S.~Doplicher and J.~E.~Roberts.
Why there is a field algebra with a compact gauge group describing the
superselection structure in particle physics.
\emph{Communications in Mathematical Physics} \textbf{131} (1990), 51--107.

\bibitem{Takesaki1970}
M.~Takesaki.
\emph{Tomita's Theory of Modular Hilbert Algebras and Its Applications}.
Lecture Notes in Mathematics \textbf{128}, Springer, 1970.

\bibitem{Connes1973}
A.~Connes.
Une classification des facteurs de type \(\mathrm{III}\).
\emph{Annales scientifiques de l'Ecole Normale Superieure} \textbf{6}
(1973), 133--252.

\bibitem{DAntoniLongo1983}
C.~D'Antoni and R.~Longo.
Interpolation by type I factors and the flip automorphism.
\emph{Journal of Functional Analysis} \textbf{51} (1983), 361--371.

\bibitem{BuchholzWichmann1986}
D.~Buchholz and E.~H. Wichmann.
Causal independence and the energy-level density of states in local quantum
field theory.
\emph{Communications in Mathematical Physics} \textbf{106} (1986), 321--344.

\bibitem{Borchers1992}
H.-J. Borchers.
The CPT theorem in two-dimensional theories of local observables.
\emph{Communications in Mathematical Physics} \textbf{143} (1992), 315--332.

\bibitem{Borchers2000}
H.-J. Borchers.
On revolutionizing quantum field theory with Tomita's modular theory.
\emph{Journal of Mathematical Physics} \textbf{41} (2000), 3604--3673.

\bibitem{Wiesbrock1993}
H.-W. Wiesbrock.
Half-sided modular inclusions of von Neumann algebras.
\emph{Communications in Mathematical Physics} \textbf{157} (1993), 83--92.

\bibitem{GlimmJaffe1987}
J.~Glimm and A.~Jaffe.
\emph{Quantum Physics: A Functional Integral Point of View}.
Second edition, Springer, 1987.

\bibitem{Hairer2014}
M.~Hairer.
A theory of regularity structures.
\emph{Inventiones Mathematicae} \textbf{198} (2014), 269--504.

\bibitem{Segal2004}
G.~Segal.
The definition of conformal field theory.
In U.~Tillmann (ed.),
\emph{Topology, Geometry and Quantum Field Theory},
London Mathematical Society Lecture Note Series \textbf{308},
Cambridge University Press, 2004, 421--577.

\bibitem{KontsevichSegal2021}
M.~Kontsevich and G.~Segal.
Wick rotation and the positivity of energy in quantum field theory.
\emph{Quarterly Journal of Mathematics} \textbf{72} (2021), 673--699.

\bibitem{Zhu1996}
Y.~Zhu.
Modular invariance of characters of vertex operator algebras.
\emph{Journal of the American Mathematical Society} \textbf{9} (1996),
237--302.

\bibitem{Huang2005}
Y.-Z. Huang.
Vertex operator algebras and the Verlinde conjecture.
\emph{Communications in Contemporary Mathematics} \textbf{10} (2008),
103--154.

\bibitem{FriedanQiuShenker1984}
D.~Friedan, Z.~Qiu, and S.~Shenker.
Conformal invariance, unitarity, and critical exponents in two dimensions.
\emph{Physical Review Letters} \textbf{52} (1984), 1575--1578.

\bibitem{GoddardKentOlive1986}
P.~Goddard, A.~Kent, and D.~Olive.
Unitary representations of the Virasoro and super-Virasoro algebras.
\emph{Communications in Mathematical Physics} \textbf{103} (1986), 105--119.

\bibitem{KawahigashiLongoMuger2001}
Y.~Kawahigashi, R.~Longo, and M.~M\"uger.
Multi-interval subfactors and modularity of representations in conformal
field theory.
\emph{Communications in Mathematical Physics} \textbf{219} (2001), 631--669.

\bibitem{CarpiKawahigashiLongoWeiner2018}
S.~Carpi, Y.~Kawahigashi, R.~Longo, and M.~Weiner.
From vertex operator algebras to conformal nets and back.
\emph{Memoirs of the American Mathematical Society} \textbf{254} (2018),
no.~1213.

\bibitem{BinghamGoldieTeugels1987}
N.~H. Bingham, C.~M. Goldie, and J.~L. Teugels.
\emph{Regular Variation}.
Encyclopedia of Mathematics and its Applications \textbf{27},
Cambridge University Press, 1987.

\bibitem{DKRV2016}
F.~David, A.~Kupiainen, R.~Rhodes, and V.~Vargas.
Liouville quantum gravity on the Riemann sphere.
\emph{Communications in Mathematical Physics} \textbf{342} (2016),
869--907.

\bibitem{KRV2020}
A.~Kupiainen, R.~Rhodes, and V.~Vargas.
Integrability of Liouville theory: proof of the DOZZ formula.
\emph{Annals of Mathematics} \textbf{191} (2020), 81--166.

\bibitem{GKRV2024}
C.~Guillarmou, A.~Kupiainen, R.~Rhodes, and V.~Vargas.
Conformal bootstrap in Liouville theory.
\emph{Acta Mathematica} \textbf{233} (2024), no.~1, 33--194.

\bibitem{Zamolodchikov1984}
Al.~B. Zamolodchikov.
Conformal symmetry in two dimensions: an explicit recurrence formula for the
conformal partial wave amplitude.
\emph{Communications in Mathematical Physics} \textbf{96} (1984), 419--422.

\bibitem{Zamolodchikov1987}
Al.~B. Zamolodchikov.
Conformal symmetry in two-dimensional space: recursion representation of
conformal block.
\emph{Theoretical and Mathematical Physics} \textbf{73} (1987), 1088--1093.

\bibitem{DO1994}
H.~Dorn and H.-J. Otto.
Two- and three-point functions in Liouville theory.
\emph{Nuclear Physics B} \textbf{429} (1994), 375--388.

\bibitem{ZZ1996}
A.~B. Zamolodchikov and Al.~B. Zamolodchikov.
Structure constants and conformal bootstrap in Liouville field theory.
\emph{Nuclear Physics B} \textbf{477} (1996), 577--605.

\bibitem{FZZ2000}
V.~Fateev, A.~Zamolodchikov, and Al.~Zamolodchikov.
Boundary Liouville field theory I: boundary state and boundary two-point
function.
\emph{arXiv:hep-th/0001012}, 2000.

\bibitem{Teschner2000}
J.~Teschner.
Remarks on Liouville theory with boundary.
\emph{Proceedings of Science} \textbf{tmr2000} (2000), 041.

\bibitem{AngRemySun2025}
M.~Ang, G.~Remy, and X.~Sun.
FZZ formula of boundary Liouville CFT via conformal welding.
\emph{Journal of the European Mathematical Society} \textbf{27} (2025),
1209--1266.

\bibitem{Seeley1967}
R.~T. Seeley.
Complex powers of an elliptic operator.
In \emph{Singular Integrals}, Proceedings of Symposia in Pure Mathematics
\textbf{10}, American Mathematical Society, 1967, 288--307.

\bibitem{Gilkey1984}
P.~B. Gilkey.
\emph{Invariance Theory, the Heat Equation, and the Atiyah--Singer Index
Theorem}.
Publish or Perish, 1984.

\bibitem{AtiyahSinger1968}
M.~F. Atiyah and I.~M. Singer.
The index of elliptic operators I.
\emph{Annals of Mathematics} \textbf{87} (1968), 484--530.

\bibitem{APS1975I}
M.~F. Atiyah, V.~K. Patodi, and I.~M. Singer.
Spectral asymmetry and Riemannian geometry I.
\emph{Mathematical Proceedings of the Cambridge Philosophical Society}
\textbf{77} (1975), 43--69.

\bibitem{APS1975II}
M.~F. Atiyah, V.~K. Patodi, and I.~M. Singer.
Spectral asymmetry and Riemannian geometry II.
\emph{Mathematical Proceedings of the Cambridge Philosophical Society}
\textbf{78} (1975), 405--432.

\bibitem{APS1976III}
M.~F. Atiyah, V.~K. Patodi, and I.~M. Singer.
Spectral asymmetry and Riemannian geometry III.
\emph{Mathematical Proceedings of the Cambridge Philosophical Society}
\textbf{79} (1976), 71--99.

\bibitem{BismutFreed1986}
J.-M. Bismut and D.~S. Freed.
The analysis of elliptic families II: Dirac operators, eta invariants, and
the holonomy theorem.
\emph{Communications in Mathematical Physics} \textbf{107} (1986), 103--163.

\bibitem{DaiFreed1994}
X.~Dai and D.~S. Freed.
Eta invariants and determinant lines.
\emph{Journal of Mathematical Physics} \textbf{35} (1994), 5155--5194.

\bibitem{Witten1982}
E.~Witten.
An \(SU(2)\) anomaly.
\emph{Physics Letters B} \textbf{117} (1982), 324--328.

\bibitem{tHooft1976Pseudoparticle}
G.~'t Hooft.
Computation of the quantum effects due to a four-dimensional pseudoparticle.
\emph{Physical Review D} \textbf{14} (1976), 3432--3450.

\bibitem{LeutwylerSmilga1992}
H.~Leutwyler and A.~Smilga.
Spectrum of Dirac operator and role of winding number in QCD.
\emph{Physical Review D} \textbf{46} (1992), 5607--5632.

\bibitem{Bernard1979InstantonDeterminants}
C.~W. Bernard.
Gauge zero modes, instanton determinants, and quantum-chromodynamic calculations.
\emph{Physical Review D} \textbf{19} (1979), 3013--3019.

\bibitem{DiakonovPetrov1986InstantonVacuum}
D.~Diakonov and V.~Yu. Petrov.
A theory of light quarks in the instanton vacuum.
\emph{Nuclear Physics B} \textbf{272} (1986), 457--489.

\bibitem{SeibergWitten1994I}
N.~Seiberg and E.~Witten.
Electric-magnetic duality, monopole condensation, and confinement in
\(\mathcal N=2\) supersymmetric Yang--Mills theory.
\emph{Nuclear Physics B} \textbf{426} (1994), 19--52.

\bibitem{SeibergWitten1994II}
N.~Seiberg and E.~Witten.
Monopoles, duality and chiral symmetry breaking in
\(\mathcal N=2\) supersymmetric QCD.
\emph{Nuclear Physics B} \textbf{431} (1994), 484--550.

\bibitem{Seiberg1994}
N.~Seiberg.
Electric-magnetic duality in supersymmetric non-Abelian gauge theories.
\emph{Nuclear Physics B} \textbf{435} (1995), 129--146.

\bibitem{MaldacenaZhiboedov2012}
J.~Maldacena and A.~Zhiboedov.
Constraining conformal field theories with a higher spin symmetry.
\emph{Journal of Physics A} \textbf{46} (2013), 214011.

\bibitem{AGY2012}
O.~Aharony, G.~Gur-Ari, and R.~Yacoby.
Correlation functions of large \(N\) Chern--Simons--matter theories and
bosonization in three dimensions.
\emph{Journal of High Energy Physics} \textbf{12} (2012), 028.

\bibitem{JainMinwallaYokoyama2013}
S.~Jain, S.~Minwalla, and S.~Yokoyama.
Chern--Simons duality with a fundamental boson and fermion.
\emph{Journal of High Energy Physics} \textbf{11} (2013), 037.

\bibitem{MTV2006BetheEigenvectors}
E.~Mukhin, V.~Tarasov, and A.~Varchenko.
Bethe eigenvectors of higher transfer matrices.
\emph{Journal of Statistical Mechanics: Theory and Experiment} \textbf{2006},
no.~8, P08002.

\bibitem{MTV2009XXXSimpleSpectrum}
E.~Mukhin, V.~Tarasov, and A.~Varchenko.
Bethe algebra of homogeneous XXX Heisenberg model has simple spectrum.
\emph{Communications in Mathematical Physics} \textbf{288} (2009), 1--42.

\bibitem{Tarasov2018Completeness}
V.~Tarasov.
Completeness of the Bethe ansatz for the periodic isotropic Heisenberg model.
\emph{Reviews in Mathematical Physics} \textbf{30} (2018), 1840018.

\bibitem{LiebRobinson1972}
E.~H. Lieb and D.~W. Robinson.
The finite group velocity of quantum spin systems.
\emph{Communications in Mathematical Physics} \textbf{28} (1972), 251--257.

\end{thebibliography}
